\section{Related Work}
\label{sec:related}

\paragraph{Global error of numerical schemes and model error in forecasting.}
The representation of a global error as a sum of transported local errors is classical in the numerical analysis of time integrators \citep[Section~II.3]{hairer1993}, where it also underlies a posteriori error estimation by adjoint (dual) problems \citep{estep1995,cao2004,giles2002}. In predictability studies, the same forced linear equation describes the growth of model error \citep{orrell2001}: \citet{nicolis2003,nicolis2004} derived the short-time evolution of the mean-square error caused by model imperfections, and \citet{vannitsem2002} showed that this evolution is linear or quadratic in time depending on the temporal correlation of the error source, the quadratic case being the generic one. The current work brings this point of view to learned surrogates, whose local errors are neither a truncation error of known order nor a prescribed stochastic process, and adds an exact decomposition of the realized error together with null models that separate the generic effect of the propagator from the structure of the defects.

\paragraph{Non-modal amplification, stochastic forcing and Lyapunov vectors.}
Non-normal linear operators can amplify perturbations transiently even when they are asymptotically stable \citep{trefethen1993,reddy1993,schmid2007}, and the optimal initial perturbation is given by the leading right singular vector of the finite-time propagator. For forced systems, generalized stability theory replaces the optimal perturbation by the response to stochastic forcing with a prescribed covariance, described by finite-time Gramians \citep{farrell1996a,farrell1996b}, and input--output analyses of fluid flows quantify the amplification from forcing to response \citep{jovanovic2005}. At long times, singular vectors of the propagator converge to Lyapunov vectors \citep{legras1996,ginelli2007,kuptsov2012}. Non-normal transient responses have also been studied in recurrent neural circuits \citep{hennequin2012} and in the optimization dynamics of neural networks \citep{ghosh2026}. Here the parameters of the surrogate are fixed, and these tools are used to characterize how the defects of a trained model populate the singular spectrum of its tangent dynamics.

\paragraph{Compounding error and long-horizon training of learned dynamics.}
The mismatch between training on reference inputs and testing on the model's own outputs is well known in sequence prediction \citep{bengio2015,venkatraman2015} and in model-based control \citep{lambert2022}. For chaotic systems, recurrent models face exploding gradients whose rate is set by the Lyapunov spectrum, which motivates teacher forcing and its generalizations \citep{mikhaeil2022,hess2023}, and statistical accuracy of learned chaotic dynamics has been related to the accuracy of their Jacobians and to shadowing \citep{park2024}. Long rollouts of neural PDE solvers and weather models are improved by multistep or unrolled objectives \citep{brandstetter2022,list2025,lam2023}, by architectural choices that control spectral content or instability \citep{lippe2023,mccabe2023}, and by inference-time corrections \citep{pedersen2025}. Jacobian regularization has been proposed to improve long-time prediction \citep{pan2018,nie2026}, and \citet{pervez2026} penalize the non-normality and non-commutativity of consecutive Jacobians, which control the transient amplification available along a rollout. These works act on the available amplification or on the training objective. The current work is diagnostic: it measures how much of the available amplification is realized by the defects of a fixed surrogate and which factor makes some forecasts worse than others, and it applies the same accounting to the effect of a multistep objective.

\paragraph{Ensembles and error estimates without reference data.}
Ensemble forecasting estimates the uncertainty of a numerical weather prediction from the spread of perturbed forecasts \citep{leutbecher2008}, and deep ensembles estimate the predictive uncertainty of neural networks from the disagreement between independently trained members \citep{lakshminarayanan2017}. A posteriori error estimates of numerical schemes, on the other hand, transport computable local residuals with the tangent or adjoint dynamics \citep{estep1995,cao2004,giles2002}. The reference-free test of the current work combines the two ideas: the disagreement between independently trained surrogates estimates the local defect, and the forced recursion transports it along the forecast.

\paragraph{Benchmarks.}
Lorenz--96 \citep{lorenz1996} and the Kuramoto--Sivashinsky equation are standard testbeds for data-driven forecasting of chaotic dynamics \citep{vlachas2018,pathak2018}. Fourier neural operators parameterize global convolutions in Fourier space \citep{li2021} and provide a surrogate class that differs from the finite-dimensional residual MLP in both representation and inductive bias.
